\chapter{Frontier Model Capabilities}

A frontier model is a model near the leading edge of capability at a given moment. The
term refers to proximity to the current frontier across one or more
dimensions: reasoning, coding, tool use, multimodality, long-context understanding,
instruction following, safety behaviour, latency.

For agent builders, frontier models matter because they determine which workflows can be
automated with acceptable reliability.

For example, for agentic coding, SWE-bench evaluates whether a frontier model can resolve real GitHub issues by modifying code
and passing the repository's own test suite.

\section{Open-Weight and Closed Models}

The model ecosystem is often described as a competition between open and closed models,
but the terminology requires precision. Closed or proprietary models are those whose
weights are withheld from public release. Users access them through an API or hosted
product, and the provider controls serving, safety behaviour, updates, and pricing.
Open-weight models are those whose weights are released, typically under a licence. Users can inspect, host,
quantise, and in many cases fine-tune them.

Closed models are often the most direct
path to frontier capability. Open-weight models offer control, data privacy, cost
optimisation, and customisation. In practice, many deployments use both: a closed model for
reasoning-intensive paths and an open-weight model for high-volume internal workflows
where cost and latency are primary constraints.


\section{The Original Transformer}

The architectural foundation worth understanding is the original Transformer, introduced
by Vaswani et al.\ in 2017. It comprised an encoder and a decoder, both built from
multi-head self-attention, position-wise feedforward networks, residual connections,
layer normalisation, and positional encodings. The encoder processed the full input
sequence with bidirectional self-attention. The decoder generated output
autoregressively, using masked self-attention over its own prior outputs and
cross-attention over the encoder's representations.

The original design targeted sequence-to-sequence tasks such as machine translation
rather than general-purpose language modelling. Its lasting contribution was the
attention mechanism as a content-dependent communication primitive: queries and keys
determine where the model attends; values carry the content aggregated at each position.
This separation of relevance from content is what allows a later token to use an earlier
definition, a function call to attend to its declaration site, or a question to locate
its answer span across a long document.

\section{The Modern Frontier Decoder}

Contemporary frontier models are almost universally decoder-only, but the internal block
has been substantially refined from the original design. The field has converged on a
representative set of architectural components, reflected across both closed models such
as GPT, Claude, and Gemini, and open-weight families such as Llama, Qwen, Mistral, and
DeepSeek.

Causal self-attention replaces the encoder-decoder cross-attention structure, with the
model attending only to prior positions during generation. Rotary Positional Embeddings
(RoPE) encode position through rotation applied to query and key vectors rather than
through additive absolute embeddings, generalising more effectively to sequence lengths
beyond those seen during training. RMSNorm, applied before each sub-layer rather than
after, provides training stability with lower computational overhead than the original
LayerNorm. Grouped-Query Attention, which has become the default for most large models,
shares key and value projections across groups of query heads, substantially reducing the
memory footprint of the KV cache during inference without significant loss in model
quality. SwiGLU, combining a Swish activation with a gated linear unit, has replaced
earlier activation functions in the feedforward sub-layer and is now the de facto
standard. FlashAttention-style kernels reorganise the attention computation to minimise
memory movement between GPU SRAM and HBM, reducing both memory usage and wall-clock
time. Sparse Mixture-of-Experts feedforward layers activate only a subset of specialised
sub-networks for each token, allowing a large total parameter count with a much smaller
active parameter count per forward pass. Multimodal encoders or adapters extend the
decoder to accept image, audio, or video inputs alongside text.

\section{How the Field Evolved}

The trajectory from the original Transformer to current agentic deployments can be read
as a sequence of distinct eras, each expanding the scope of what language models could
do.

The representation era was defined by encoder-only models such as BERT, which produced
strong contextual embeddings through bidirectional attention and achieved state-of-the-art
results on classification, extraction, and ranking tasks. These models established the
value of deep contextual representations but were designed for understanding rather than
generation.

The text-to-text era unified diverse tasks under a single sequence-to-sequence interface.
Models such as T5 reframed translation, summarisation, question answering, and
classification as instances of the same text-in, text-out paradigm, demonstrating that
a single model architecture and training objective could cover a broad range of
downstream applications.

The autoregressive scaling era established that next-token prediction on massive corpora,
combined with scaling of data, parameters, and compute, could produce broad and
transferable capabilities. The GPT family demonstrated this trajectory, and the empirical
scaling laws that emerged from this period shaped research investment across the field.

The instruction-following era introduced supervised fine-tuning and preference
optimisation, including reinforcement learning from human feedback, to align base
models with human instructions. The model's task shifted from text continuation to
goal-directed response generation, enabling deployment as a general-purpose assistant.

The reasoning era introduced inference-time compute scaling: models trained with
verifiable rewards in mathematics, code, and formal settings learned to allocate
additional computation to difficult problems, producing longer and more deliberate
reasoning chains before committing to an answer.

The agentic era optimises models to call tools reliably, maintain coherent state across
long contexts, and operate within loops that incorporate external feedback from tool
results and environment observations. This is the setting for which the rest of this
book provides architectural guidance.

This progression explains why agentic AI is a continuation of the language modelling
research programme rather than a departure from it. Agents are the systems layer built
on top of the modelling stack; their capabilities and limitations are directly shaped by
where the underlying model sits on this trajectory.

\section{Summary}

Frontier models are the capability substrate of agentic systems. The original Transformer
introduced attention as a content-dependent communication mechanism; modern frontier
decoders extend that foundation with long-context positional representations, efficient
attention variants, normalisation improvements, sparse mixture-of-experts layers,
post-training alignment, reasoning optimisation, and multimodal input handling. Model
selection should be grounded in a capability profile evaluated inside the actual agent
workflow, with benchmark results interpreted as system-level measurements rather than
intrinsic model properties.
